\chapter{Route to Corollary 3.12}
\label{ch:corollary-route}

\section{Diagnostic finite label and Gaussian layer}

\begin{definition}[Finite label square-weight profile]
  \label{def:finite-label-profile}
  \lean{Iut.Stage1.IUTStage1ZModSquareWeightProfile,Iut.Stage1.IUTStage1ZModSquareWeightProfile.GaussianMonoidDegreeEvaluation}
  \leanok
  The current finite model records the \(j^2\)-weighted Gaussian degree
  behavior of theta labels.  It separates raw coordinate square weights,
  sign-quotient full labels, zero and nonzero classes, and several averaged
  log-volume conventions.
\end{definition}

\begin{proposition}[Diagnostic: no silent scalar absorption]
  \label{prop:no-silent-scalar-absorption}
  \uses{def:finite-label-profile}
  \lean{Iut.Stage1.Experiments.no_labelIndependent_transport_scale_absorbs_j2}
  \leanok
  At the representative level, one label-independent scalar cannot absorb all
  \(j^2\)-weighted theta degrees when the environment degree is nonzero.  This
  is a counter-collapse diagnostic: it shows that the formalization does not
  silently erase the \(q^{j^2}\) issue inside a real-line identification, but
  it does not discharge the source Step (xi) hull/log-volume comparison.
\end{proposition}

\begin{proof}
  This is represented in the finite-label experiment surface.  The exact Lean
  declarations distinguish representative square weights from descended
  full-label data and from the final hull/log-volume endpoint.
\end{proof}

\section{Nonarchimedean upper-semi route}

\begin{definition}[Local \(\Ind{3}\) orientation]
  \label{def:ind3-local-orientation}
  \lean{Iut.Stage1.Experiments.Ind3FirstPassDashboard,Iut.Stage1.Experiments.ind3FirstPassDashboard}
  \leanok
  The experiment dashboard is a Lean-checked status report.  It records that
  missing real alignment blocks the final route, that nonarchimedean
  \(\Ind{3}\) entries have the current packet-to-theta orientation, and that
  archimedean entries have the reverse orientation for this route.  It should
  not be read as a source construction of the \(\Ind{3}\) comparison.
\end{definition}

\begin{theorem}[Nonarchimedean entry reaches the signed comparison]
  \label{thm:nonarchimedean-entry-final-q-theta}
  \uses{def:finite-label-profile,def:ind3-local-orientation,def:stage1-comparison-endpoint}
  \lean{Iut.Stage1.Experiments.nonarchimedeanEntry_finalQTheta}
  \leanok
  A nonarchimedean local \(\Ind{3}\) entry, together with the required
  real-line/log-volume alignments and structured \(\SHE\) square-weight
  transport audit, implies
  \[
    q_{\mathrm{signed}}\le \Theta_{\mathrm{signed}}.
  \]
\end{theorem}

\begin{proof}
  The theorem composes the local nonarchimedean alignment with the
  Hodge-descent packet-to-theta route.  It still assumes the relevant
  source-facing records; it proves that those records have the correct formal
  shape to feed the endpoint.
\end{proof}

\begin{theorem}[Strongest current packet-calibrated route]
  \label{thm:packet-calibrated-route}
  \uses{thm:nonarchimedean-entry-final-q-theta,thm:strict-boundary-alternatives}
  \lean{Iut.Stage1.IUTStage1SourcePackage.PlaceAuditedMultiradialThetaHullEndpoint.LogVolumeChartAudit.FLZModCuspLabelThetaHodgeDescentPacketTransportAudit.boundarySignedEqualityOrStrictCTheta_of_gaussianKummerForgettingPacketCalibratedEntry}
  \leanok
  In the strongest current route, Gaussian identity at the canonical full
  label \(1\), the finite Kummer-plus-forgetting product-log-volume chain,
  packet-normalized theta-source equality, packet-normalized \(\Ind{3}\)
  target calibration, and packet-normalized entry-target calibration imply
  \[
    \bigl(q_{\mathrm{signed}}=\Theta_{\mathrm{signed}}
       \;\wedge\; \Theta_{\mathrm{signed}}<0\bigr)
       \quad\text{or}\quad
    -1<\CTheta.
  \]
\end{theorem}

\begin{proof}
  The route first constructs the needed nonarchimedean log-Kummer upper-semi
  compatibility from the Kummer transfer and mono-analytic forgetting chain.
  It then derives the target-side equalities through packet-normalized
  calibration rather than taking the direct theta-to-entry equality as a
  primitive input.
\end{proof}

\begin{definition}[Packet-calibrated hypothesis ledger]
  \label{def:packet-calibrated-hypothesis-ledger}
  \uses{thm:packet-calibrated-route}
  \notready
  The strongest current route consumes a large family of source-facing
  hypotheses: a source package, hull/determinant obligations, a log-volume
  chart audit, a Hodge-descent transport packet, source and target Gaussian
  profiles and evaluations, canonical full-label \(1\) preservation, source
  and target log-volume equalities, finite Kummer transfer, mono-analytic
  forgetting, a nonarchimedean entry and membership proof, packet-normalized
  calibrations, \qPilot{} positivity, and the displayed \(\CTheta\) upper
  bound.  This node is intentionally marked not ready: the theorem checks
  the conditional route after these inputs are supplied, but the source
  construction of the inputs remains the mathematical work.
\end{definition}

\section{The Corollary 3.12 node}

\begin{definition}[\(\Theta\) finite-real endpoint]
  \label{def:theta-finite-real-endpoint}
  \lean{Iut.Stage1.Experiments.corollary312StatementEndpoint_theta_ne_plusInfinity}
  \notready
  Corollary 3.12 starts from an extended signed \(\Theta\) log-volume and must
  conclude that the endpoint is finite real-valued.  The current Lean theorem
  projects this fact from a statement-endpoint record; the source construction
  of that record remains open.
\end{definition}

\begin{definition}[\qPilot{} positivity]
  \label{def:q-pilot-positivity}
  \lean{Iut.Stage1.Experiments.corollary312StatementEndpoint_q_neg}
  \notready
  The final \(\CTheta\) algebra requires \(0<-q_{\mathrm{signed}}\), or
  equivalently that the signed \qPilot{} log-volume is negative.  The current
  theorem extracts this from the statement-endpoint record; a source proof
  must still construct the record from the prime-strip data.
\end{definition}

\begin{definition}[Displayed \(\CTheta\) upper bound]
  \label{def:displayed-ctheta-upper-bound}
  \lean{Iut.Stage1.Experiments.corollary312StatementEndpoint_realTheta_le_cTheta_absLogQ}
  \notready
  The displayed source hypothesis
  \[
    \Theta_{\mathrm{signed}}
      \le \CTheta\cdot(-q_{\mathrm{signed}})
  \]
  is a separate input to the ordered-real lower bound for \(\CTheta\).  It
  should remain visible rather than being folded into the signed comparison.
\end{definition}

\begin{theorem}[\(\CTheta\) lower bound from signed data]
  \label{thm:ctheta-lower-bound-from-signed-data}
  \uses{def:q-pilot-positivity,def:displayed-ctheta-upper-bound,thm:final-real-calculation}
  \lean{Iut.Stage1.Experiments.corollary312SignedCTheta_lower_bound}
  \leanok
  Once the signed comparison, \qPilot{} positivity, and displayed
  \(\CTheta\) upper bound have been supplied, the ordered-real conclusion
  \(-1\le\CTheta\) is Lean-checked algebra.
\end{theorem}

\begin{proof}
  This is not source mathematics; it is the real-inequality consequence of the
  signed endpoint data.
\end{proof}

\begin{corollary}[Conditional Corollary 3.12 endpoint]
  \label{cor:conditional-corollary-312}
  \uses{prop:step-xi-upper-ray-input,lem:ledger-three-term-comparison,lem:signed-comparison-gives-corollary,def:theta-finite-real-endpoint,def:q-pilot-positivity,def:displayed-ctheta-upper-bound,thm:ctheta-lower-bound-from-signed-data,thm:packet-calibrated-route,def:packet-calibrated-hypothesis-ledger}
  \lean{Iut.Stage1.SourceObligationLedger.corollary312}
  \notready
  If the source constructions of \IUT{} III supply the required
  multiradial output, \(\IPL\), \(\SHE\), \(\APT\), hull-plus-determinant
  comparison, log-Kummer upper-semi compatibility, and packet-normalized
  calibrations, then the Corollary 3.12-shaped signed pilot inequality follows.
  The Lean tag points to the conditional route theorem from the source ledger,
  not merely to the bare predicate that states the target inequality.
\end{corollary}

\begin{proof}
  The Lean code already proves the endpoint from the named route records.
  What is not yet proved is that the source papers construct all of those
  records without collapsing the histories, labels, or comparison levels that
  the dispute requires us to keep separate.
\end{proof}
